\begin{abstract}
Tokenization serves as a fixed interface between raw symbolic inputs and neural encoders, yet most existing tokenizers are optimized for large language models and natural language distributions.
When used as drop-in components in structured and semi-structured prediction pipelines, such tokenizers induce uncontrolled compression that distorts role-dependent semantics and wastes token budget under tight compute constraints.
We formalize tokenization as a constrained interface design problem and introduce \emph{directed semantic distortion}, a prefix-aligned measure of irrecoverable task-relevant information loss induced by the interface.
Based on this formulation, we propose \textbf{Controlled Interface Tokenization (CIT)}, which enforces a versioned interface contract with role-explicit boundaries and typed-value integrity, and constructs task-specialized vocabularies by maximizing compression efficiency within the contract under deterministic execution.
Directed semantic distortion is used to evaluate and calibrate semantic fidelity across tokenizers and budgets, yielding empirical rate--distortion frontiers.
Across synthetic and public structured classification tasks, CIT achieves superior rate--distortion trade-offs, higher accuracy under matched token budgets, and improved robustness to semantics-preserving re-serialization and prefix-limited prediction.
\end{abstract}

\section{Introduction}
\label{sec:intro}

Tokenization is often treated as a benign preprocessing step.
In encoder-only systems operating under hard sequence and latency budgets, it is instead a fixed \emph{representation interface}: inputs are deterministically serialized into strings and then irrevocably segmented into tokens, after which all computation is conditioned on the token sequence.
In this strict drop-in setting, changing the tokenizer changes the interface contract between raw symbolic inputs and the encoder, and therefore changes which distinctions can \emph{ever} be represented, learned, and recovered under binding token budgets.

Structured and semi-structured streams make this interface problem acute.
Their semantics is role- and boundary-dependent: the same character substring may be a key, a delimiter, or a value depending on its position and surrounding syntax.
They also contain high-cardinality fields (identifiers, timestamps, hashes, URLs) whose internal character patterns are unstable but whose span integrity is task-critical.
Under deterministic left-to-right, prefix-emitting execution, tokenization induces a causal mapping from raw prefixes to tokenized prefixes.
We call the resulting failure mode \emph{prefix semantic collapse}: two raw prefixes that correspond to different task-relevant semantic states become indistinguishable because they map to the same observable token prefix.
Collapse is an interface property: once it happens, no downstream encoder operating under the same tokenizer can reconstruct a distinction that never reaches the model, regardless of model size or training recipe.
This is especially damaging for streaming or early-decision settings and for any deployment with hard maximum sequence lengths.

Canonical tokenizers fail on structured inputs in a structural way.
WordPiece and BPE optimize subword statistics largely driven by natural language, without constraints that respect role boundaries or typed value spans.
As a result, they may segment across boundary-bearing symbols, entangle role cues with incidental value fragments, and become brittle to semantics-preserving re-serialization (e.g., benign field reordering, whitespace/escaping normalization).
Conversely, using finer-grained pieces to avoid brittleness often leads to heavy-tailed token lengths on high-cardinality fields, inflating attention cost and truncation risk.

These interface effects appear immediately under standard encoder limits.
In a scan of six widely used pretrained tokenizers on 5{,}000 CSIC 2010 HTTP requests with $L_{\max}{=}512$, average sequence lengths range from 359--468 tokens and 7--29\% of examples exceed the limit (29.4\% for BERT WordPiece; 7.3\% for a GPT-3.5/GPT-4-family tokenizer).
The tail is heavy: the longest request reaches 1{,}034 tokens, illustrating that truncation risk is an interface property of segmentation rather than a modeling detail.

% \begin{figure}[t]
% \centering
% \begin{tokviz}
% \tok{GET}\tok{\detokenize{ }}\tok{/}\tok{search}\tok{?}\tok{q}\tok{=}\tok{\detokenize{%}}\tok{E7}\tok{\detokenize{%}}\tok{82}\tok{\detokenize{%}}\tok{B8}\tok{\detokenize{%}}\tok{E9}\tok{\detokenize{%}}\tok{B8}\tok{\detokenize{%}}\tok{A1}%
% \tok{\&}\tok{token}\tok{=}\tok{a}\tok{B}\tok{1}\tok{!}\tok{x}\tok{Y}\tok{9}\tok{@}\tok{z}\tok{Z}\tok{2}\tok{\#}%
% \tok{\&}\tok{id}\tok{=}\tok{0}\tok{x}\tok{7}\tok{f}\tok{8}\tok{e}\tok{9}\tok{a}\tok{0}\tok{b}\tok{2}\tok{c}\tok{4}\tok{d}\tok{\detokenize{ }}\tok{HTTP}\tok{/}\tok{1.1}\\
% \tok{Host}\tok{:}\tok{\detokenize{ }}\tok{127.0.0.1}\tok{:}\tok{8080}\\
% \tok{X-Trace-ID}\tok{:}\tok{\detokenize{ }}\tok{550e8400}\tok{-}\tok{e29b}\tok{-}\tok{41d4}\tok{-}\tok{a716}\tok{-}\tok{446655440000}\\
% \tok{Cookie}\tok{:}\tok{\detokenize{ }}\tok{sess}\tok{=}\tok{q}\tok{Wer}\tok{Ty}\tok{Ui}\tok{Op}\tok{12345}
% \end{tokviz}
% \vspace{-0.3em}
% \caption{Token-boundary visualization for a representative semi-structured request under a canonical WordPiece tokenizer (BERT family). Colors indicate token boundaries. The segmentation fragments boundary-bearing symbols and high-cardinality fields into many short pieces, obscuring role/value structure and inflating effective token budget.}
% \label{fig:intro_token_example}
% \vspace{-0.8em}
% \end{figure}

We therefore treat tokenization for structured prediction as a constrained interface design problem.
We introduce \emph{directed semantic distortion}, a prefix-aligned measure of irrecoverable task-relevant information loss incurred when prediction must be made from tokenized prefixes rather than raw prefixes.
This yields a deployment-aligned rate--distortion view of tokenizers under vocabulary and sequence constraints.
Building on this formulation, we propose \textbf{Controlled Interface Tokenization (CIT)}: a deterministic tokenizer that (i) enforces an explicit, versioned interface contract with role-preserving boundaries and typed-value integrity, and (ii) constructs task-specialized vocabularies by maximizing compression gain within the contract-feasible space under deployed prefix-emitting execution.

\paragraph{Contributions.}
Our contributions are:
\begin{itemize}[leftmargin=1.25em, itemsep=0.2em, topsep=0.2em]
    \item We identify \emph{prefix semantic collapse} as an interface failure mode for structured streams: distinct task-relevant prefix states become indistinguishable after prefix-emitting tokenization, creating irrecoverable information loss under any downstream encoder.
    \item We propose \emph{directed semantic distortion}, a prefix-aligned, task-relevant measure of interface-induced information loss, enabling empirical rate--distortion analysis for tokenization interfaces.
    \item We introduce \textbf{CIT}, a contract-constrained, deterministic tokenization system that preserves role boundaries and typed-value integrity while inducing deployment-aligned vocabularies under prefix-emitting execution.
    \item We empirically validate a causal chain from interface to learning to system behavior: CIT avoids collapse on a synthetic prefix stress test, improves end-to-end learning and drift robustness under token-fair training on a public HTTP benchmark, and shifts the accuracy--latency Pareto frontier under fully trained encoder models.
\end{itemize}

The remainder of the paper develops the interface formulation, presents CIT, and evaluates its efficiency, robustness, and end-to-end impact under strict drop-in constraints.


\section{Problem Formulation}
\label{sec:formulation}

We formalize tokenization as a representation interface for structured and semi-structured prediction under a strict drop-in regime.
Our goal is to characterize how interface design induces an irreversible compression of task-relevant information, and to make this loss explicit and comparable across tokenizers.

\subsection{Setup and Drop-in Regime}
\label{sec:setup}

Each input example originates from a structured or semi-structured record $r \in \mathcal{R}$, such as a tabular row, a key--value map, or a protocol/log event.
To enable the use of encoder-only neural models, records are deterministically serialized into symbolic strings
\begin{equation}
X \;=\; s(r) \in \Sigma^\ast,
\label{eq:serialization}
\end{equation}
where $\Sigma$ is a finite alphabet of characters or bytes.
This record-to-text practice is widely used to reuse sequence models for non-linguistic inputs, including tabular and semi-structured prediction pipelines~\cite{yin2020tabert,herzig2020tapas}.

A tokenizer is a mapping $\tau : \Sigma^\ast \rightarrow \mathcal{V}^\ast$ over a finite vocabulary $\mathcal{V}$, producing a token sequence
\[
Z \;=\; \tau(X).
\]
We focus on a strict \emph{drop-in regime}, in which the downstream encoder architecture, optimization procedure, and training budget are fixed (e.g., a given Transformer encoder~\cite{vaswani2017attention,devlin2019bert}).
Under this regime, the tokenizer is the primary modifiable component, and its execution must be deterministic, reproducible, and compatible with deployment constraints.

We assume \emph{prefix-emitting} execution: tokens are emitted left-to-right and never revised.
For a raw prefix $X_{1:t}$, the observable tokenized prefix is
\begin{equation}
Z^{(t)} \;=\; \tau(X_{1:t}),
\label{eq:prefix_tok}
\end{equation}
which is irrevocably determined by the prefix.
This assumption models streaming and early-decision settings and induces a causal interface between the raw input stream and the model's available observations at inference time.

\subsection{Semantic Prefix States and Interface Irreversibility}
\label{sec:semantic_state}

Prediction on structured inputs is inherently prefix-dependent.
Let $Y$ denote the task semantics (e.g., a class label), and let $\ell$ be the task loss (e.g., cross-entropy).
For each raw prefix $X_{1:t}$, we define the \emph{semantic prefix state} as the Bayes-optimal conditional distribution
\begin{equation}
\mathcal{S}(X_{1:t})
\;\triangleq\;
P^\star(\,\cdot \mid X_{1:t}),
\label{eq:semantic_raw}
\end{equation}
which captures the strongest achievable semantics given direct access to the raw prefix.
This definition abstracts away nuisance variation and retains exactly the information relevant for prediction under the chosen loss~\cite{cover1999elements,tishby2000information}.

Tokenization restricts access to the raw input by exposing only the tokenized prefix $Z^{(t)}$.
Fix a reference predictor family $\mathcal{Q}$ whose members map tokenized prefixes to predictive distributions over $Y$.
The best achievable semantics under the interface at prefix length $t$ is
\begin{equation}
\mathcal{S}_\tau(Z^{(t)})
\;\triangleq\;
\arg\min_{q \in \mathcal{Q}}
\;
\mathbb{E}\!\left[\ell\!\left(q(\,\cdot \mid Z^{(t)}), Y\right)\right].
\label{eq:semantic_tok}
\end{equation}

Because $\tau$ is deterministic and prefix-emitting, the mapping $X_{1:t} \mapsto Z^{(t)}$ is generally many-to-one.
Distinct raw prefixes can therefore induce the same observable token prefix, forcing any token-observing predictor to average over their semantics.
This identifies a fundamental \emph{irreversibility} at the interface: once two raw prefixes collapse to the same $Z^{(t)}$, no downstream encoder-only model can distinguish them without modifying the interface itself.
In structured settings, such collapse commonly arises when segmentation crosses role boundaries or fragments high-cardinality values whose internal characters are semantically irrelevant but whose integrity is role-critical~\cite{yin2020tabert,herzig2020tapas}.

\subsection{Directed Semantic Distortion}
\label{sec:distortion}

We quantify irrecoverable semantic loss induced by tokenization through a prefix-aligned, directed distortion.
Intuitively, distortion measures how much task-relevant semantic resolution is lost when prediction must be made from $Z^{(t)}$ rather than directly from $X_{1:t}$.

Let $\mathcal{R}^\star(t)$ denote the Bayes risk at prefix length $t$,
\begin{equation}
\mathcal{R}^\star(t)
\;\triangleq\;
\mathbb{E}\!\left[\ell\!\left(\mathcal{S}(X_{1:t}), Y\right)\right],
\label{eq:bayes_raw}
\end{equation}
and let $\mathcal{R}_\tau(t)$ denote the minimum achievable risk when only the tokenized prefix is observed,
\begin{equation}
\mathcal{R}_\tau(t)
\;\triangleq\;
\min_{q \in \mathcal{Q}}
\mathbb{E}\!\left[\ell\!\left(q(\,\cdot \mid Z^{(t)}), Y\right)\right].
\label{eq:bayes_tok}
\end{equation}
The \emph{directed semantic distortion} of tokenizer $\tau$ is defined as the average excess risk over prefixes:
\begin{equation}
\Delta(\tau)
\;\triangleq\;
\mathbb{E}\!\left[
\frac{1}{T}\sum_{t=1}^{T}
\big(\mathcal{R}_\tau(t) - \mathcal{R}^\star(t)\big)
\right].
\label{eq:distortion}
\end{equation}

This definition makes three properties explicit.
First, $\Delta(\tau)\ge 0$ by construction: restricting observations from $X_{1:t}$ to $Z^{(t)}$ cannot improve the best achievable risk.
Second, $\Delta(\tau)$ is an \emph{interface-level} quantity: it depends only on what information the tokenizer preserves about prefixes and upper-bounds the performance of any downstream encoder-only model operating under the same interface.
Third, prefix averaging penalizes interfaces that delay task-relevant information to later tokens, which is critical under truncation and token-budget constraints.

In practice, $\Delta(\tau)$ is not directly observable because it depends on Bayes-optimal semantics.
In our experiments, we estimate a surrogate distortion using a teacher--student construction, where a reference predictor observes raw prefixes and a restricted predictor observes tokenized prefixes.
We use this surrogate \emph{only for evaluation and calibration}, so that measured gaps reflect interface-induced information loss rather than representation adaptation~\cite{hinton2015distilling}.

\subsection{Rate--Distortion Trade-off under Drop-in Constraints}
\label{sec:rate_distortion}

Tokenization induces a compression of the raw input prior to learning.
We quantify interface cost by the expected token length
\begin{equation}
R(\tau) \;\triangleq\; \mathbb{E}[|Z|],
\label{eq:rate}
\end{equation}
which directly controls effective context length and inference cost in encoder-only models.

Let $B$ denote a vocabulary budget, and define the drop-in feasible set
\begin{equation}
\mathcal{T}_{\mathrm{drop}}(B)
\;=\;
\big\{
\tau \;\big|\;
|\mathcal{V}| \le B,\;
\tau \text{ deterministic and prefix-emitting}
\big\}.
\label{eq:feasible}
\end{equation}
Each tokenizer $\tau \in \mathcal{T}_{\mathrm{drop}}(B)$ induces a pair $(R(\tau), \Delta(\tau))$, representing its interface efficiency and irrecoverable semantic loss.

The resulting \emph{rate--semantic distortion frontier}
\begin{equation}
\mathcal{F}(B)
\;\triangleq\;
\big\{
(R(\tau), \Delta(\tau)) :
\tau \in \mathcal{T}_{\mathrm{drop}}(B)
\big\}
\label{eq:frontier}
\end{equation}
characterizes the fundamental trade-off imposed by the interface under a fixed vocabulary budget.
Lowering $R(\tau)$ increases usable context and reduces compute, while excessive compression necessarily increases $\Delta(\tau)$ by collapsing task-relevant distinctions at the prefix level.
This mirrors classical rate--distortion trade-offs in lossy compression~\cite{shannon1948mathematical,cover1999elements}, with the crucial difference that distortion here is task-aligned and prefix-sensitive rather than reconstruction-based.

Rather than treating tokenization as a heuristic preprocessing choice, we view interface design as selecting an operating point on $\mathcal{F}(B)$.
The next section develops a deployment-aligned method for constructing tokenizers that improve this trade-off under deterministic, prefix-emitting execution semantics.

\section{Method}
\label{sec:method}

We now present \emph{Controlled Interface Tokenization (CIT)}, a deployment-aligned tokenization system designed to improve the rate--distortion trade-off under strict drop-in constraints.
CIT operationalizes the interface perspective developed in Section~\ref{sec:formulation} by combining an explicit interface contract, a deterministic execution model, and a vocabulary construction procedure aligned with deployed behavior.

\subsection{Interface Contract}
\label{sec:contract}

CIT is defined with respect to an explicit \emph{interface contract} that specifies which segmentations are semantically admissible.
The contract constrains the tokenization space to prevent interface-induced semantic collapse before any vocabulary induction is performed.

\paragraph{Role-explicit boundaries.}
Structured and semi-structured records encode semantics through explicit roles, such as field names, delimiters, and separators.
The interface contract declares a set of role boundaries that tokenization is not permitted to cross.
This prevents merges that conflate distinct semantic roles, such as key--value separators or field terminators.

\paragraph{Typed-value integrity.}
Many structured domains contain high-cardinality values (e.g., identifiers, timestamps, addresses) whose internal character structure is not semantically meaningful but whose integrity is role-critical.
The contract specifies a set of typed-value classes and enforces their atomicity: typed values are treated as indivisible units, and neither segmentation nor vocabulary induction may split them.
This constraint prevents interface collapse caused by fragmenting unstable value instances.

Together, these constraints define a contract-feasible space of tokenizations that preserve role semantics and value integrity by construction.
All tokenizers evaluated under CIT share the same contract, ensuring that comparisons isolate the effect of vocabulary construction rather than upstream preprocessing.

\subsection{Deterministic Prefix-Emitting Execution}
\label{sec:execution}

CIT adopts a deterministic, prefix-emitting execution model.
Given a contract-transformed input string, token emission proceeds left-to-right and never revises past decisions.
At each position, the tokenizer emits the longest contract-feasible substring present in the vocabulary; if none exists, a fallback atomic unit (character or byte) is emitted.

This execution model satisfies two deployment-critical properties.
First, it ensures that tokenized prefixes depend only on observed raw prefixes, aligning execution with the causal interface assumed in Section~\ref{sec:formulation}.
Second, determinism guarantees reproducibility and allows quantities estimated during vocabulary construction to reflect true runtime behavior.
Ties are resolved deterministically (e.g., by token identifier order), yielding a unique segmentation for each input.

\subsection{Vocabulary Construction via Deployment-Aligned Compression Gain}
\label{sec:construction}

Starting from an atomic contract vocabulary $\mathcal{V}_0$, CIT constructs a vocabulary of size $B$ by iteratively adding contract-feasible tokens.
Let $\tau_{\mathcal{V}}$ denote the tokenizer induced by vocabulary $\mathcal{V}$ under the deterministic execution policy described above.

We quantify interface cost by the expected token length $R(\tau)=\mathbb{E}[|Z|]$.
For a candidate token $c$ admissible under the interface contract, we define its \emph{compression gain} as
\begin{equation}
G(c;\mathcal{V})
\;\triangleq\;
R(\tau_{\mathcal{V}}) - R(\tau_{\mathcal{V}\cup\{c\}}),
\label{eq:gain}
\end{equation}
that is, the reduction in expected token length when $c$ is added to the vocabulary and executed under deployed semantics.

Unlike frequency-based criteria used in classical subword induction, $G(c;\mathcal{V})$ is evaluated under the same deterministic, prefix-emitting execution model used at inference.
As a result, it directly reflects interface-level efficiency improvements rather than abstract substring statistics.

\subsection{Greedy Construction under Interface Constraints}
\label{sec:greedy}

CIT expands the vocabulary by repeatedly selecting the contract-feasible token with the largest compression gain:
\begin{equation}
c_k
\;\in\;
\arg\max_{c \in \mathcal{C}(\mathcal{V}_k)}
G(c;\mathcal{V}_k),
\label{eq:greedy}
\end{equation}
where $\mathcal{C}(\mathcal{V}_k)$ denotes the set of contract-feasible candidates under the current vocabulary.
The vocabulary is updated as $\mathcal{V}_{k+1} \leftarrow \mathcal{V}_k \cup \{c_k\}$ until the budget $B$ is reached.

Directed semantic distortion is \emph{not} optimized inside the construction loop.
Instead, CIT separates \emph{construction} from \emph{evaluation}: vocabulary induction is driven by deployment-aligned compression gain, while semantic fidelity is assessed post hoc using the distortion metric introduced in Section~\ref{sec:distortion}.
This separation yields a predictable and reproducible construction procedure while enabling empirical rate--distortion analysis across tokenizers and budgets.

\subsection{Execution Alignment and Deployment Properties}
\label{sec:properties}

A defining feature of CIT is \emph{execution alignment}: the same deterministic, prefix-emitting segmentation semantics are used during vocabulary construction and at deployment.
This alignment ensures that marginal quantities estimated during construction correspond to actual runtime behavior.

\begin{lemma}[Prefix invariance]
\label{lem:prefix}
Let $\tau_{\mathcal{V}}$ be the tokenizer induced by vocabulary $\mathcal{V}$ under deterministic, prefix-emitting execution.
For any input string $X$ and prefix length $t$, the tokenized prefix $\tau_{\mathcal{V}}(X_{1:t})$ is independent of future characters $X_{t+1:}$.
\end{lemma}

Lemma~\ref{lem:prefix} follows directly from left-to-right, non-revising token emission and guarantees that tokenized prefixes observed during construction are identical to those observed at deployment.
As a result, interface cost estimates and downstream evaluations are free from training--deployment mismatch.

Finally, once the vocabulary $\mathcal{V}_B$ is fixed, the tokenizer $\tau_{\mathcal{V}_B}$ can be compiled into a deterministic finite-state matcher.
This compilation eliminates dynamic control flow, guarantees linear-time execution, and produces a standalone artifact suitable for latency-sensitive systems.
\section{Background}
\label{sec:background}

\subsection{Interface setting and drop-in constraints}
\label{sec:bg_setting}

We study tokenization as a fixed representation interface for supervised prediction on structured and
semi-structured symbolic inputs.
Each example originates from a structured record $r \in \mathcal{R}$, such as a tabular row, a key--value map,
or a protocol/log event.
To enable the use of encoder-only neural encoders, records are deterministically serialized into symbolic
strings
\begin{equation}
X \;=\; s(r) \in \Sigma^\ast,
\label{eq:bg_serialization}
\end{equation}
where $\Sigma$ is an alphabet of characters or bytes.
This record-to-text practice is widely adopted to reuse sequence models for non-linguistic inputs, including
tabular and semi-structured prediction pipelines~\cite{yin2020tabert,herzig2020tapas}.

A tokenizer is a mapping $\tau : \Sigma^\ast \rightarrow \mathcal{V}^\ast$ over a finite vocabulary
$\mathcal{V}$, producing a token sequence $Z = \tau(X)$.
We focus on a strict \emph{drop-in regime} in which the downstream encoder architecture, optimization procedure,
and training budget are fixed (e.g., a given Transformer encoder)~\cite{vaswani2017attention,devlin2019bert}.
Under this regime, the tokenizer is the primary modifiable component, and its execution must be deterministic,
reproducible, and compatible with deployment constraints.

We assume \emph{prefix-emitting} execution: tokens are emitted left-to-right and never revised.
For a raw prefix $X_{1:t}$, the observable tokenized prefix is
\begin{equation}
Z^{(t)} \;=\; \tau(X_{1:t}),
\label{eq:bg_prefix}
\end{equation}
which is irrevocably determined by the prefix.
This requirement models streaming and early-decision settings, and it induces a causal interface between the
raw input stream and the model's available observations at inference time.

In encoder-only Transformers, inference cost scales at least linearly—and often superlinearly—with token
sequence length due to self-attention~\cite{vaswani2017attention,kaplan2020scaling}.
Consequently, under small-model and low-latency constraints, tokenization directly controls effective context
length, truncation behavior, and compute allocation.
This makes the tokenization interface a first-order bottleneck for both efficiency and accuracy in structured
prediction systems.

\subsection{Semantic prefix states and irreversibility of tokenization}
\label{sec:bg_semantics}

Prediction on structured inputs is inherently prefix-dependent.
Let $Y$ denote task semantics (e.g., a class label), and let $\ell$ be the task loss (e.g., cross-entropy).
For each raw prefix $X_{1:t}$, define the \emph{semantic prefix state} as the Bayes-optimal conditional
distribution
\begin{equation}
\mathcal{S}(X_{1:t})
\;\triangleq\;
P^\star(\,\cdot \mid X_{1:t}),
\label{eq:semantic_state_raw}
\end{equation}
which captures the strongest achievable semantics given direct access to the raw prefix.
This notion parallels classical views of task-relevant information: $X_{1:t}$ contains a mixture of
predictive signal and nuisance variation, and the semantic state represents exactly what matters for prediction
under the chosen loss~\cite{tishby2000information,cover1999elements}.

Tokenization restricts access to the raw input by exposing only the tokenized prefix
$Z^{(t)} = \tau(X_{1:t})$.
To formalize what is achievable through the interface, fix a reference predictor family $\mathcal{Q}$ whose
members map tokenized prefixes to predictive distributions over $Y$.
Define the best token-observing semantics at time $t$ as
\begin{equation}
q^\star_\tau(\,\cdot \mid Z^{(t)})
\;\in\;
\arg\min_{q \in \mathcal{Q}}
\;
\mathbb{E}\!\left[\ell\!\left(q(\,\cdot \mid Z^{(t)}), Y\right)\right],
\label{eq:semantic_state_token}
\end{equation}
and let $\mathcal{S}_\tau(Z^{(t)}) \triangleq q^\star_\tau(\,\cdot \mid Z^{(t)})$.

Because $\tau$ is deterministic and prefix-emitting, the mapping $X_{1:t} \mapsto Z^{(t)}$ is generally many-to-one.
Distinct raw prefixes can therefore induce the same observable token prefix, forcing any token-observing
predictor to average over their semantics.
This identifies an \emph{irreversibility} at the interface: once two raw prefixes collapse to the same
$Z^{(t)}$, no downstream encoder-only model can distinguish them without changing the interface.
In structured settings, this collapse can occur when segmentation crosses role boundaries or fragments
high-cardinality values whose internal characters are not semantically meaningful but whose integrity is
role-critical~\cite{yin2020tabert,herzig2020tapas}.
The interface-design problem is thus not merely compression, but compression that preserves task-relevant
distinctions needed for prediction from prefixes.

\subsection{Directed semantic distortion}
\label{sec:bg_distortion}

We quantify irrecoverable semantic loss induced by tokenization through a prefix-aligned, directed distortion.
Intuitively, distortion measures how much semantic resolution is lost when one must predict from
$Z^{(t)}$ rather than directly from $X_{1:t}$.

Let $\mathcal{R}^\star(t)$ denote the Bayes risk at prefix length $t$,
\begin{equation}
\mathcal{R}^\star(t)
\;\triangleq\;
\mathbb{E}\!\left[\ell\!\left(\mathcal{S}(X_{1:t}), Y\right)\right],
\label{eq:bayes_risk_raw}
\end{equation}
and let $\mathcal{R}_\tau(t)$ denote the best achievable risk given only the tokenized prefix,
\begin{equation}
\mathcal{R}_\tau(t)
\;\triangleq\;
\min_{q \in \mathcal{Q}}
\mathbb{E}\!\left[\ell\!\left(q(\,\cdot \mid Z^{(t)}), Y\right)\right].
\label{eq:bayes_risk_tok}
\end{equation}
The directed (raw $\rightarrow$ tokenized) semantic distortion of $\tau$ is defined as the average excess risk
over prefix lengths:
\begin{equation}
\Delta(\tau)
\;\triangleq\;
\mathbb{E}\!\left[
\frac{1}{T}\sum_{t=1}^{T}
\big(\mathcal{R}_\tau(t) - \mathcal{R}^\star(t)\big)
\right].
\label{eq:semantic_distortion}
\end{equation}

This definition makes three aspects explicit.
First, $\Delta(\tau)\ge 0$ by construction: restricting observations from $X_{1:t}$ to $Z^{(t)}$ cannot improve
the best achievable risk under the same loss.
Second, $\Delta(\tau)$ is an \emph{interface property}: it depends on what information $\tau$ preserves about
prefixes, and it upper-bounds the performance one can obtain from any downstream encoder-only model without
changing the interface.
Third, the prefix averaging penalizes interfaces that delay task-relevant information to later tokens, which is
critical under truncation and token-budget constraints.


Operationally, $\Delta(\tau)$ can be estimated using a teacher--student construction: a teacher captures reference semantics from raw prefixes, while a student is restricted to tokenized prefixes. In our experiments, we use lightweight surrogate predictors to estimate this gap \emph{for evaluation}, avoiding full-scale training so that the metric reflects interface-induced information loss rather than representation adaptation. This viewpoint aligns with the use of KL or cross-entropy gaps as measures of irrecoverable information loss in distillation-style settings~\cite{hinton2015distilling} and connects to information-theoretic perspectives on task-relevant compression~\cite{tishby2000information,cover1999elements}.

\subsection{Rate--distortion frontier under drop-in constraints}
\label{sec:bg_frontier}

Tokenization induces an interface-level compression of the raw input prior to learning.
We quantify interface cost by the expected token length
\begin{equation}
R(\tau) \;\triangleq\; \mathbb{E}[|Z|],
\label{eq:rate_def}
\end{equation}
which directly controls effective context length and inference cost in encoder-only models.

Let $B$ denote a vocabulary budget, and define the drop-in feasible set
\begin{equation}
\mathcal{T}_{\mathrm{drop}}(B)
\;=\;
\big\{
\tau \;\big|\;
|\mathcal{V}| \le B,\;
\tau \text{ deterministic and prefix-emitting}
\big\}.
\label{eq:dropin_set}
\end{equation}
Each feasible tokenizer $\tau \in \mathcal{T}_{\mathrm{drop}}(B)$ induces a pair $(R(\tau), \Delta(\tau))$,
representing its efficiency and irrecoverable semantic loss.

We define the \emph{rate--semantic distortion frontier}
\begin{equation}
\mathcal{F}(B)
\;\triangleq\;
\big\{
(R(\tau), \Delta(\tau)) :
\tau \in \mathcal{T}_{\mathrm{drop}}(B)
\big\},
\label{eq:frontier}
\end{equation}
which characterizes the fundamental trade-off imposed by the interface under a fixed vocabulary budget.
Lowering $R(\tau)$ increases usable context and reduces compute, while excessive compression necessarily
increases $\Delta(\tau)$ by collapsing task-relevant distinctions at the prefix level.
This mirrors classical rate--distortion trade-offs in lossy compression~\cite{shannon1948mathematical,cover1999elements,blahut2003computation,arimoto1972algorithm},
with the crucial difference that our distortion is task-aligned and prefix-sensitive rather than based on
reconstruction fidelity.

Rather than treating tokenization as a heuristic preprocessing choice, we view interface design as selecting an
operating point on $\mathcal{F}(B)$.
A standard scalarization of this trade-off is given by the Lagrangian objective
\begin{equation}
\min_{\tau \in \mathcal{T}_{\mathrm{drop}}(B)}
\;\;
R(\tau) + \lambda\,\Delta(\tau),
\label{eq:lagrangian_objective}
\end{equation}
where $\lambda \ge 0$ controls the relative importance of semantic fidelity.
Directly optimizing~\eqref{eq:lagrangian_objective} is intractable due to the combinatorial nature of the
feasible set and the dependence of $\Delta(\tau)$ on unobservable Bayes semantics.
The next section develops a deployment-aligned relaxation and an algorithm that approximately optimizes this
trade-off under deterministic, prefix-emitting execution semantics.

\section{Algorithms}
\label{sec:algorithms}

This section develops a concrete procedure for constructing tokenization interfaces that improve the rate--semantic distortion trade-off under strict drop-in constraints.
A central challenge is that the distortion term depends on Bayes-optimal semantics and is not directly computable.
Rather than attempting to optimize distortion during vocabulary induction (which would be prohibitively expensive and sensitive to surrogate choices), we separate \emph{construction} from \emph{calibration}:
CIT constructs a deterministic, contract-feasible tokenizer by maximizing compression efficiency within a constrained feasible set, and then uses surrogate directed semantic distortion to evaluate semantic fidelity and characterize empirical rate--distortion frontiers.
This design yields a deployment-aligned algorithm with predictable runtime and reproducible artifacts.


\subsection{Deployment-aligned distortion for interface evaluation}
\label{sec:alg_relaxation}

The directed semantic distortion $\Delta(\tau)$ depends on Bayes-optimal semantic states and is therefore not directly computable.
We introduce a deployment-aligned surrogate that replaces Bayes semantics with restricted-capacity predictors trained on observable data.
Let $\widehat{\mathcal{S}}(X_{1:t})$ and $\widehat{\mathcal{S}}_\tau(Z^{(t)})$ denote semantic states induced by such predictors when observing raw prefixes and tokenized prefixes, respectively.
These predictors are intentionally lightweight so that differences reflect information preserved by the interface rather than model expressiveness.

We define the surrogate semantic distortion
\begin{equation}
\widehat{\Delta}(\tau)
\;\triangleq\;
\mathbb{E}\!\left[
\frac{1}{T}
\sum_{t=1}^{T}
\mathcal{L}\big(
\widehat{\mathcal{S}}(X_{1:t}),
\widehat{\mathcal{S}}_\tau(Z^{(t)})
\big)
\right],
\label{eq:surrogate_distortion}
\end{equation}
where $\mathcal{L}$ is a discrepancy measure such as cross-entropy or KL divergence.
In CIT, $\widehat{\Delta}(\tau)$ is used to \emph{evaluate and calibrate} semantic fidelity across tokenizers and vocabulary budgets, enabling empirical rate--distortion analysis under the same deterministic, prefix-emitting execution semantics used at deployment.


\subsection{Contract-feasible candidates and compression gain}
\label{sec:alg_marginal}

CIT constructs tokenizers within a constrained feasible set induced by an explicit interface contract.
The contract specifies (i) role-explicit boundary markers, and (ii) typed-value integrity constraints for high-cardinality fields (e.g., identifiers, timestamps, addresses), preventing merges that cross semantic roles or fragment values in semantically invalid ways.
Let $\mathcal{V}_0$ denote the atomic contract vocabulary.

Given a current vocabulary $\mathcal{V}$ and the induced deterministic, prefix-emitting tokenizer $\tau_{\mathcal{V}}$, consider adding a contract-feasible candidate token $c$ that respects all boundary and typing constraints.
We quantify interface cost by expected token length $R(\tau)=\mathbb{E}[|Z|]$.
The \emph{compression gain} of adding $c$ is defined as the reduction in expected token length under deployed execution:
\begin{equation}
G(c;\mathcal{V})
\;\triangleq\;
R(\tau_{\mathcal{V}}) - R(\tau_{\mathcal{V}\cup\{c\}}).
\label{eq:gain_def}
\end{equation}
Unlike raw substring frequency used in classical subword induction, $G(c;\mathcal{V})$ is evaluated under the same prefix-emitting semantics used at deployment and respects contract constraints.

\subsection{Greedy construction under interface constraints}
\label{sec:alg_greedy}

Starting from $\mathcal{V}_0$, CIT expands the vocabulary by repeatedly selecting the contract-feasible token that maximizes compression gain:
\begin{equation}
c_k
\;\in\;
\arg\max_{c \in \mathcal{C}(\mathcal{V}_k)}
G(c;\mathcal{V}_k),
\label{eq:alg_selection_gain}
\end{equation}
and updating $\mathcal{V}_{k+1}\leftarrow \mathcal{V}_k\cup\{c_k\}$ until $|\mathcal{V}_k|=B$.
This yields a deterministic and deployment-aligned vocabulary construction procedure with predictable runtime.

Directed semantic distortion is \emph{not} optimized inside the main construction loop.
Instead, we compute $\widehat{\Delta}(\tau)$ offline to validate semantic fidelity and to map empirical rate--distortion operating points across tokenizers and budgets.
Optionally, one may refine the greedy choice by re-ranking a small top-$K$ shortlist of candidates using $\widehat{\Delta}$; we use this only as an implementation detail when needed and report $\widehat{\Delta}$ primarily as an evaluation metric.

\begin{algorithm}[t]
\caption{Controlled Interface Tokenization (CIT)}
\label{alg:cit}
\begin{algorithmic}[1]
\STATE \textbf{Input:} Training corpus $\mathcal{D}$, vocabulary budget $B$, interface contract $\mathcal{K}$
\STATE \textbf{Initialize:} $\mathcal{V}_0 \leftarrow$ atomic contract vocabulary induced by $\mathcal{K}$
\FOR{$k = 0$ \textbf{to} $B-|\mathcal{V}_0|-1$}
    \STATE Construct contract-feasible candidate set $\mathcal{C}(\mathcal{V}_k)$
    \FOR{each $c \in \mathcal{C}(\mathcal{V}_k)$}
        \STATE Estimate compression gain $G(c;\mathcal{V}_k)$ under deployed execution (Eq.~\eqref{eq:gain_def})
    \ENDFOR
    \STATE Select $c_k$ using Eq.~\eqref{eq:alg_selection_gain}
    \STATE Update $\mathcal{V}_{k+1} \leftarrow \mathcal{V}_k \cup \{c_k\}$
\ENDFOR
\STATE \textbf{Output:} Vocabulary $\mathcal{V}_B$ and tokenizer $\tau_{\mathcal{V}_B}$
\end{algorithmic}
\end{algorithm}


\subsection{Execution alignment and deployment properties}
\label{sec:alg_execution}

A defining feature of CIT is execution alignment: the same deterministic, prefix-emitting segmentation semantics
are used during vocabulary construction and at inference.
This ensures that marginal quantities estimated during optimization reflect true runtime behavior.

\begin{lemma}[Prefix invariance]
\label{lem:prefix_invariance}
Let $\tau_{\mathcal{V}}$ be the tokenizer induced by vocabulary $\mathcal{V}$ under deterministic,
prefix-emitting execution.
For any input string $X$ and prefix length $t$, the tokenized prefix $\tau_{\mathcal{V}}(X_{1:t})$ is
independent of future characters $X_{t+1:}$.
\end{lemma}

\begin{proof}
Token emission proceeds left-to-right and never revises past decisions.
Therefore, at position $t$, emitted tokens depend only on $X_{1:t}$ and $\mathcal{V}$.
\end{proof}

Lemma~\ref{lem:prefix_invariance} implies that tokenized prefixes observed during construction are identical to
those observed at deployment.
As a result, $\widehat{\mathcal{J}}_\lambda$ is a deployment-aligned objective rather than a training-time
artifact.

Finally, once the vocabulary $\mathcal{V}_B$ is fixed, the tokenizer $\tau_{\mathcal{V}_B}$ can be compiled into
a deterministic finite-state transducer.
This compilation eliminates dynamic control flow, guarantees linear-time execution, and produces a standalone
artifact suitable for latency-sensitive systems.

\subsection{Interface properties and algorithmic guarantees}
\label{sec:alg_properties}

We now formalize two properties that clarify the role of the tokenization interface and the behavior of
the greedy construction procedure.
The first establishes that information loss induced by the interface is fundamentally irrecoverable by
any downstream encoder-only model.
The second shows that the proposed greedy algorithm is consistent with the deployment-aligned objective
introduced in Section~\ref{sec:bg_frontier}.

\begin{proposition}[Irrecoverability of interface-induced prefix collapse]
\label{prop:irrecoverability}
Let $\tau$ be a deterministic, prefix-emitting tokenizer, and let
$X_{1:t}$ and $X'_{1:t}$ be two raw prefixes such that
\[
\tau(X_{1:t}) = \tau(X'_{1:t}) = Z^{(t)}.
\]
Then, for any token-observing predictor $q \in \mathcal{Q}$ and any task loss $\ell$,
the predictive distribution conditioned on the interface must coincide:
\[
q(\,\cdot \mid Z^{(t)}) \text{ is identical for both prefixes.}
\]
Consequently, if the Bayes-optimal semantic states differ,
$\mathcal{S}(X_{1:t}) \neq \mathcal{S}(X'_{1:t})$,
then the excess Bayes risk induced by the interface at prefix length $t$ is strictly positive.
\end{proposition}

Proposition~\ref{prop:irrecoverability} makes explicit that semantic distinctions collapsed by the
tokenization interface cannot be recovered by any downstream encoder-only model without modifying the
interface itself.
This property underlies the definition of directed semantic distortion in
Section~\ref{sec:bg_distortion} and motivates optimizing tokenization as an interface-level design problem
rather than relying on downstream model capacity.

\begin{proposition}[Monotonic descent under deployment-aligned greedy construction]
\label{prop:greedy_descent}
Let $\widehat{\mathcal{J}}_\lambda$ denote the relaxed interface objective defined in
Eq.~\eqref{eq:relaxed_objective}, and let
$\{\mathcal{V}_k\}_{k \ge 0}$ be the sequence of vocabularies produced by
Algorithm~\ref{alg:cit} under fixed deterministic, prefix-emitting execution semantics.
Then the relaxed objective is non-increasing along the construction path:
\[
\widehat{\mathcal{J}}_\lambda(\tau_{\mathcal{V}_{k+1}})
\;\le\;
\widehat{\mathcal{J}}_\lambda(\tau_{\mathcal{V}_{k}})
\quad \text{for all } k.
\]
\end{proposition}

Proposition~\ref{prop:greedy_descent} establishes that the greedy selection rule used by CIT performs a
consistent local descent on the deployment-aligned objective.
While global optimality is not guaranteed due to non-submodularity and interaction effects between tokens,
this property distinguishes CIT from frequency-driven heuristics that lack an explicit interface-level
objective.

\section{Experiments}
\label{sec:experiments}

We re-anchor the empirical evaluation around three effect experiments that trace a causal chain from interface behavior to learning dynamics to system-level outcomes:
\textbf{E1} asks whether prefix-emitting tokenization induces semantic collapse and whether CIT prevents it;
\textbf{E2} asks whether preventing collapse makes end-to-end learning more sample-efficient and stable under a fixed token budget;
\textbf{E3} asks whether these interface gains translate into better accuracy--latency trade-offs under fully trained encoder models.
Across experiments, we keep the serialization schema, model architecture family, and optimizer settings fixed, varying only the tokenization interface.

\subsection{Experimental protocol}
\label{sec:exp_protocol}

\paragraph{Datasets and serialization.}
All inputs originate from structured or semi-structured records and are deterministically
serialized into symbolic strings using a fixed schema shared across all tokenizers.
We report results on (i) a synthetic structured prefix task designed to stress role/boundary dependence and semantics-preserving drift,
(ii) CSIC 2010 HTTP anomaly detection (public, structured, and drift-prone),
and (iii) OpenML Adult for an end-to-end accuracy/latency Pareto slice.
(Additional public tabular benchmarks and vocabulary sweeps are deferred to the appendix.)

\paragraph{Tokenizers.}
We compare CIT against standard frequency-driven tokenizers:
BPE~\cite{sennrich2016neural}, WordPiece~\cite{schuster2012japanese}, and
Unigram~\cite{kudo2018subword}, all trained on the same data and constrained to the same
vocabulary budget ($|\mathcal{V}|{=}2048$).
All tokenizers are deterministic and prefix-emitting.

\paragraph{Training budgets and models.}
We match training compute using token-fair budgets: each run processes the same number of non-padding encoder tokens under a shared maximum sequence length and truncation policy.
E1 uses a tiny encoder and a 3M-token budget ($L_{\max}{=}128$) to isolate interface failure modes.
E2 trains end-to-end Tiny/Small encoders on CSIC under a 10M-token budget ($L_{\max}{=}512$), and also reports a step-fair control with matched optimizer steps.
E3 trains end-to-end encoder families on Adult under a 5M-token budget ($L_{\max}{=}512$) and measures inference latency.

\paragraph{Metrics.}
We report average token length (\emph{rate}), tail token length (95th percentile),
surrogate directed semantic distortion $\widehat{\Delta}(\tau)$,
classification accuracy, robustness under semantics-preserving re-serialization drift,
and, where applicable, end-to-end inference latency (95th percentile).

\subsection{E1: Does CIT Prevent Prefix Semantic Collapse?}
\label{sec:exp_e1}

E1 isolates interface correctness in a controlled synthetic setting.
We construct structured records whose labels depend on a role-specific signal placed mid-record, while other fields contain high-cardinality values and boundary-bearing markers.
We evaluate (i) surrogate directed semantic distortion (prefix-level information loss) and (ii) robustness under semantics-preserving re-serializations (field shuffling, whitespace/number-format normalization, and insertion of irrelevant dummy fields).

Semantic collapse is not hypothetical: with BPE/WordPiece/Unigram, test accuracy collapses to near chance on this task (4.5--5.6\% for 20-way classification), consistent with role-dependent prefixes becoming indistinguishable under the interface.
In contrast, CIT prevents collapse and attains 100\% test accuracy.
The interface diagnostics agree: CIT reduces surrogate distortion from 0.67--0.75 (baselines) to 0.295, and compresses the token rate from 645--692 tokens/example (baselines) to 81 tokens/example.
Under semantics-preserving drift, baselines remain near chance (3.4--5.3\% drift accuracy), while CIT retains 85.4\% accuracy.

This experiment validates directed semantic distortion as an interface-level metric:
tokenizers with lower $\widehat{\Delta}(\tau)$ exhibit markedly improved drift robustness, even under compute-matched training.

\subsection{E2: Token-Fair End-to-End Learning on CSIC 2010}
\label{sec:exp_e2}

E2 tests the learning benefit of avoiding collapse under a strict compute budget.
We train Tiny/Small encoder models end-to-end on CSIC 2010 HTTP anomaly detection under a fixed token-fair budget of 10M encoder tokens.
We also evaluate drift robustness under realistic, semantics-preserving re-serialization: reordering record-aware key--value groups and normalizing whitespace.

Under token-fair training, CIT is substantially more efficient (217 avg.\ tokens/example, P95 416) than BPE/WordPiece (\(\sim\)327--329 avg., P95 514--526) and Unigram (406 avg., P95 595).
This efficiency translates into more parameter updates under the same token budget (e.g., 1425 steps for CIT vs.\ \(\sim\)950 for BPE/WordPiece in our setup), and faster convergence in tokens: for Tiny models, CIT reaches 0.88 validation accuracy after 2.4M tokens, whereas BPE/WordPiece require \(\sim\)9M tokens to reach the same level.
In final test accuracy, CIT achieves 91.2\% on Tiny models (vs.\ 89.0\% WordPiece, 85.3\% BPE), and remains competitive on Small models (90.1\% vs.\ 90.3--90.5\% for BPE/WordPiece).

Robustness reveals the interface effect most clearly: under drift, CIT retains its performance (drift accuracy equals in-distribution accuracy), while BPE/WordPiece suffer catastrophic drops (e.g., 49.6\% and 46.9\% drift accuracy for Tiny; 47.9\% and 51.1\% for Small).
Unigram is comparatively robust but substantially less efficient and less accurate under the same token budget.
A step-fair control with matched optimizer steps preserves CIT's drift robustness while attenuating accuracy differences, consistent with robustness arising from interface correctness and accuracy gains under token-fair training arising from better token efficiency.

\subsection{E3: End-to-End Accuracy--Latency Trade-offs}
\label{sec:exp_e3}

E3 closes the loop at the system level by measuring whether CIT shifts the end-to-end accuracy/latency trade-off under fully trained models.
On OpenML Adult, we train each (tokenizer, model family) pair end-to-end under a fixed token budget (5M encoder tokens) and measure average/P95 sequence lengths and P95 forward-pass latency.

CIT reduces average sequence length from \(\sim\)106 tokens (BPE/WordPiece) to 60.5 tokens (P95 69), yielding a consistent latency reduction.
For the Mini model family, CIT attains 0.849 accuracy with 1.15ms P95 latency, compared to 0.843/1.35ms for BPE and 0.850/1.36ms for WordPiece.
For the Small family, CIT matches WordPiece accuracy (0.761) while reducing P95 latency from 3.64ms to 2.52ms.
Overall, CIT provides improved operating points on the accuracy--latency Pareto frontier by simultaneously reducing interface cost and preserving task-relevant prefix information (lower surrogate distortion than both baselines).

\subsection{Summary}
\label{sec:exp_summary}

Across E1--E3, we observe a consistent chain: semantic collapse is real under prefix-emitting interfaces; CIT prevents collapse by enforcing role- and type-aware interface constraints; this improves end-to-end learning under token-fair budgets and yields better accuracy--latency trade-offs under fully trained encoder models.
Additional vocabulary sweeps and public tabular benchmarks are provided in the appendix.


\section{Discussion}
\label{sec:discussion}

Our results support the view that tokenization for domain-specific prediction should prioritize specialization over generality.
On structured inputs, frequency-driven compression alone is insufficient and can actively harm task performance by conflating semantic roles.
CIT demonstrates that enforcing interface constraints and aligning vocabulary induction with task semantics yields more efficient and robust representations, particularly for small encoder-only models.


\section{Conclusion}
\label{sec:conclusion}

We formulated tokenization as controlled interface compression under a drop-in regime and introduced directed semantic distortion as a principled interface metric.
We proposed Controlled Interface Tokenization, aligning feasible constraints, vocabulary induction, and deterministic execution.
Across synthetic and public structured benchmarks, CIT improves rate--distortion trade-offs, downstream accuracy under matched training budgets, and robustness to re-serialization and prefix-limited prediction.

\section*{Impact Statement}

This work improves efficiency and robustness in structured prediction pipelines under fixed compute budgets.
Applications should follow responsible data handling and privacy-preserving practices.
